\section{机器人、VLA 与预训练下半场}

访谈后半段多次回到 robotics。谢赛宁的判断是：在谈 AGI 或 super intelligence 之前，先问能否做出足够可靠、能在家庭环境里承担家务的机器人。很多几岁或十几岁的孩子能做的事，今天的机器人还做不好。这不是四肢硬件单点问题，而是机器人大脑问题。

\subsection{VLA 与世界模型的分工}

VLA 即 Vision-Language-Action，目标是把视觉、语言指令和动作连接起来。它是机器人路线的重要形态，但谢赛宁认为，仅靠把 language model 当 foundation，然后加 action head，不足以解决 world model pre-training 的问题。VLA 可以在具体任务上很强，却未必承担“预训练下半场”的基础层工作。

\begin{knowledgebox}{术语消化：机器人路线中的几个关键词}
\begin{tabularx}{\textwidth}{p{0.22\textwidth}p{0.32\textwidth}X}
\toprule
术语 & 解决的问题 & 本期中的含义 \\
\midrule
VLA & 视觉、语言和动作的端到端连接 & 机器人执行任务的重要路线，但不等同于完整世界模型。 \\
Robot Brain & 机器人的大脑或上游智能层 & 需要感知、记忆、预测、规划和控制。 \\
Hardware Scaling Law & 通过更多机器人部署获得更多真实数据和硬件经验 & 本体公司必须面对，但不直接解决大脑预训练。 \\
Imitation Learning & 从示范轨迹学习动作策略 & 短期可用，但依赖数据和任务分布。 \\
World Model Pre-training & 面向连续多模态信号的基础预训练 & 谢赛宁称为预训练的下半场。 \\
\bottomrule
\end{tabularx}
\end{knowledgebox}

\subsection{预训练下半场输入什么、输出什么}

主持人追问 world model pre-training 输入什么、输出什么。谢赛宁的回答是：至少长期看，输入应是连续空间、高维、有噪声的多模态信号，开始可以是 video，也可能包括除视觉外的其他 encoder。输出是什么仍是 research question。这个“不知道”很重要：它说明问题还没被过早收敛成固定 recipe。

\begin{importantbox}{预训练下半场的开放性}
语言模型预训练有清晰形式：next-token prediction。世界模型预训练还没有同样稳定的范式。可能的输入是视频、传感器和多模态流；可能的目标是预测、表征、惊讶、行动后果或可规划状态。这里仍是基础研究问题。
\end{importantbox}

\subsection{本章小结}

机器人是世界模型最自然的出口之一，但机器人公司的短期资源常被硬件部署和具体任务牵引。世界模型路线要补的是上游大脑：能理解物理世界、保留记忆、预测后果，并为 VLA 和机器人控制提供底座。

